\documentclass[../../main.tex]{subfiles}% 注意这里的文件路径不能用 ./main.tex ,否则用latexmk编译子文件会报错
\graphicspath{{\subfix{./image/}}} % 指定图片目录,后续可以直接使用图片文件名
% 注意这里的文件路径不能用 ../../image/ ,否则用latexmk编译子文件会报错

% 例如:
% \begin{figure}[H]
% \centering
% \includegraphics[scale=0.3]{图.png}
% \caption{}
% \label{figure:图}
% \end{figure}
% 注意:上述\label{}一定要放在\caption{}之后,否则引用图片序号会只会显示??.

\begin{document}

\section{正则参数曲线}

\begin{definition}[参数曲线]
设区间$[a,b]\subset\mathbb{R}$，将连续映射
\[
\bm{r}: [a,b] \to \mathbb{R}^3,\quad t\mapsto \bm{r}(t)=\big(x(t),y(t),z(t)\big)
\]
称为\textbf{参数曲线}，还把参数$t$增大的方向称为该参数曲线的\textbf{正向}.

若$\bm{r}(t)$还满足$x(t),y(t),z(t)$均为连续可微函数，则称$\bm{r}(t)$为$\mathbb{R}^3$中的\textbf{连续可微参数曲线}，称$t$为\textbf{参数}。
\end{definition}

\begin{definition}[切向量与正则点]
设$\bm{r}(t)$是连续可微参数曲线，根据导数的定义，将
\[
\bm{r}'(t_0) =\lim_{\Delta t\rightarrow 0} \frac{r\left( t_0+\Delta t \right) -r\left( t \right)}{\Delta t} = \left(x'(t_0),y'(t_0),z'(t_0)\right).
\]
称为参数曲线$\bm{r}(t)$在$t_0$处的\textbf{切向量}。

若$\bm{r}'(t_0)\neq \bm{0}$,则称$t_0$为\textbf{正则点}。若$\bm{r}'(t)=\bm{0}$,则称$t_0$为\textbf{非正则点}。

\end{definition}

\begin{proposition}
设$r(t)$是连续可微参数曲线，$t_0$是$r(t)$的一个正则点，则
\begin{enumerate}[(1)]
\item $r'(t_0)$的方向正好指向曲线$r(t)$的正向.

\item 曲线在正则点$t_0$的切线$\boldsymbol{X}\left( u \right)$的方程是
\begin{align*}
\boldsymbol{X}\left( u \right) =r\left( t_0 \right) +ur\left( t_0 \right) ,\quad \forall u\in \mathbb{R} .
\end{align*}
\end{enumerate}
\end{proposition}
\begin{proof}
定义自明.

\end{proof}

\begin{definition}[正则参数曲线和参数变换]
若连续可微参数曲线$\bm{r}(t)$满足
\begin{enumerate}[(1)]
\item $\bm{r}(t)$至少是自变量$t$的三次以上连续可微的向量函数；

\item 处处是正则点，即$\bm{r}'(t)\neq \bm{0}$，$\forall t\in[a,b]$.
\end{enumerate}
则称曲线$\bm{r}(t)$为\textbf{正则参数曲线}。

设$t=t(u)$是$\mathbb{R}$上的变换，且满足
\begin{enumerate}[(1)]
\item $t=t(u)$是$\mathbb{R}$上的双射;

\item $t(u)$是$u$的三次以上连续可微函数;

\item $t'(u)$处处不为零.
\end{enumerate}
则称$t=t(u)$是\textbf{参数变换}.
\end{definition}
\begin{remark}
要求$\bm{r}(t)$至少三次以上连续可微因为：曲线的几何不变量涉及$\bm{r}(t)$的三次导数。
\end{remark}

\begin{theorem}
设$\bm{r_1}(u),\bm{r_2}(u)$是正则参数曲线，记关系$\sim_1,\,\sim_2$分别为
\begin{gather*}
\boldsymbol{r}_1(u)\sim_1 \boldsymbol{r}_2(u)\Longleftrightarrow \text{存在参数变换}t=t\left( u \right) ,\text{使}\boldsymbol{r}_1\left( u \right) =\boldsymbol{r}_2\left( t\left( u \right) \right),
\\
\boldsymbol{r}_1(u)\sim_2 \boldsymbol{r}_2(u)\Longleftrightarrow \text{存在参数变换}t=t\left( u \right) ,\text{满足}t\prime \left( u \right) >0\left( \forall u\in \mathbb{R} \right) ,\text{使}\boldsymbol{r}_1\left( u \right) =\boldsymbol{r}_2\left( t\left( u \right) \right) .
\end{gather*}
则关系$\sim_1,\,\sim_2$都是等价关系.由关系$\sim_1$确定的一个等价类称为\textbf{一条正则曲线}.
由关系$\sim_2$确定的一个等价类称为\textbf{一条有向的正则曲线}.
\end{theorem}
\begin{note}
关系$\sim_1$中的参数变换保持曲线的定向不变.
\end{note}
\begin{proof}
只需注意到
\begin{align*}
r^{\prime}\left( u \right) =\frac{\mathrm{d}}{\mathrm{d}u}r\left( t\left( u \right) \right) =r^{\prime}\left( t\left( u \right) \right) \cdot t^{\prime}\left( u \right) ,
\end{align*}
再结合定义立得.

\end{proof}







1. \textbf{显式表示}：$y=y(x),\ z=z(x)$，取$x$为参数，$\bm{r}(x)=\big(x,y(x),z(x)\big)$，必为正则曲线（因$\bm{r}'(x)=(1,y'(x),z'(x))\neq\bm{0}$）。

2. **隐式表示**：由方程组
\[
\begin{cases}
f(x,y,z)=0,\\
g(x,y,z)=0
\end{cases}
\]
确定，若Jacobi矩阵
\[
\begin{pmatrix}
\dfrac{\partial f}{\partial x}&\dfrac{\partial f}{\partial y}&\dfrac{\partial f}{\partial z}\\
\dfrac{\partial g}{\partial x}&\dfrac{\partial g}{\partial y}&\dfrac{\partial g}{\partial z}
\end{pmatrix}
\]
的秩为2，则在交点邻域内确定一条正则曲线。

\begin{example}
半三次曲线$\bm{r}(t)=(t^2,t^3)$，$t=0$为非正则点，切线方向在该点发生$180^\circ$翻转。
\end{example}

\begin{exercise}
\begin{enumerate}
\item 半径为$r$的圆盘沿$x$轴无滑滚动，求盘上距中心$a(0<a\leq r)$的点的轨迹方程（摆线）。

\item 求曲线$\bm{r}(t)=(2\sin t, t, \cos t)$在$t=0$和$t=\frac{\pi}{2}$处的切线方程。

\item 设$p$为曲线$\bm{r}(t)$外定点，$\rho(t)=|\bm{r}(t)-p|$，若$\rho(t)$在$t_0$处取极值，证明$\bm{r}'(t_0)\perp (\bm{r}(t_0)-p)$。
\end{enumerate}
\end{exercise}











\end{document}